\documentclass[11pt]{article}
\usepackage[margin=0.87in]{geometry}
\usepackage[T1]{fontenc}
\usepackage{lmodern,microtype,amsmath,amssymb,booktabs,tabularx,graphicx}
\usepackage{xcolor}
\usepackage[colorlinks=true,linkcolor=blue!45!black,urlcolor=blue!50!black]{hyperref}
\usepackage{enumitem}
\setlist{nosep,leftmargin=*}
\setlength{\parskip}{0.4em}
\setlength{\parindent}{0pt}
\setlength{\emergencystretch}{1em}
\graphicspath{{images/}}
\input{values.tex}
\title{Stock evolution, local release, and the remaining matching problem}
\author{Research calculation for the Entropic Scalar EFT project}
\date{23 October 2026}
\begin{document}
\maketitle
\begin{abstract}
A low-density constitutive completion of the capacity-clock action preserves
the tested homogeneous and linear cosmological branch while supplying a
positive pressure response before local exhaustion. Its total energy has a
constant lower bound, its regular future-directed current remains nonnegative,
and its material characteristics are subluminal. The construction uses the
existing clock and an auxiliary density; its constitutive function and energy
floor are additional effective assumptions. Independent calculations constrain
the other proposed connections. Passive relaxation toward a declining
equilibrium has the wrong sign to suppress an initially equilibrated high-redshift
galaxy response. The existing two-pole quantum model obeys a noise--delay
inequality. Relabeling stock and its small cosmological mass loss do not select
the galactic equilibrium. The minimal Einstein action supplies two luminal
tensor modes, while the native continuum and the physical contact operator
remain open. Forty-two standalone Lean certificates check the stated algebra,
calculus and comparison lemmas. They do not certify an observational likelihood
or a complete microscopic theory.
\end{abstract}

\section{Result and scope}
The five questions from the empirical assessment have different closure
requirements. Table~\ref{tab:status} records the outcome of this calculation.
The strongest new construction concerns local release. It avoids changing the
already-tested linear cosmology by modifying the action only in a separate
low-density regime. The stock and response calculations identify a common
missing physical ingredient: the interaction that changes the orbital
distribution and selects the galactic response.

\begin{table}[h]
\small
\begin{tabularx}{\textwidth}{@{}p{0.18\textwidth}XX@{}}
\toprule
Connection & Established here & Remaining physical input \\
\midrule
Stock evolution & Conservative kinetic/action framework; conversion and
adiabatic-scaling restrictions & The interaction selecting the RAR equilibrium,
assembly and director dynamics \\
Response regimes & Lag sign; quantum noise--delay bound; rate nonuniqueness &
Carrier assignment, initialization, relaxation and spatial correlation scales \\
Local release & Selected positive constitutive completion with unchanged
linear branch & Selection of its constitutive law; nonlinear spatial and
multistream validation \\
Tensor sector & Two luminal modes for the specified minimal effective action &
Native continuum, operator matching and allowed curvature couplings \\
Matching & Explicit families preserving the relevant identities while changing
predictions & A physical matrix element or independent observable selecting the
remaining coefficients \\
\bottomrule
\end{tabularx}
\caption{Effective closure, formal proof and physical selection are recorded separately.}
\label{tab:status}
\end{table}

The current manuscripts and upstream Lean files were treated as inputs. The
GitHub head checked for this work was
\href{https://github.com/jakechinitz/Entropic-EFT-Lean/commit/432da8c4a8b169a17146b53e4155d6b37915a12f}{\texttt{432da8c4}}.
The local supplement was checked with the available Lean/mathlib 4.19 runtime;
the full upstream 4.35 library was not rebuilt. Existing numerical results are
identified as such. New calculations use the fitted exponential release
history from the unified empirical program.

\section{A positive local completion that preserves the linear branch}
\subsection{The action and the depleted regime}
Use signature $(-,+,+,+)$ and $X=-\nabla_\mu\phi\nabla^\mu\phi/2$.
The original constrained-clock matter Lagrangian is
\begin{equation}
 P_{\rm dust}=\lambda(X-\tfrac12)-V(\phi).
\end{equation}
On its unit-clock branch, $\lambda$ is the pressureless capacity density.
A prescribed positive $V_\phi$ continues to drain a sufficiently empty
element through $\lambda=0$. The integrating-factor positivity lemma does
not apply when its fractional rate is $Q/\lambda$ and becomes singular.

Choose a constant $K$ larger than $V$ on the interval of interest, and define
\begin{equation}
 L(\phi)=3[K-V(\phi)]>0,\qquad \lambda\ge L(\phi)/2.
\end{equation}
Consider the covariant action
\begin{equation}\label{eq:action}
 S=\int d^4x\sqrt{-g}\left[\frac{M_{\rm Pl}^2}{2}\mathcal R
       +\lambda X-e(\lambda,\phi)-V(\phi)\right]+S_b[g],
\end{equation}
where
\begin{equation}\label{eq:e}
 e(\lambda,\phi)=\frac{\lambda}{2}
       +\frac{2}{3L(\phi)^2}[L(\phi)-\lambda]_+^3.
\end{equation}
The auxiliary variable has no derivative term. The cubic joins the dust
branch with equal value, first derivative and second derivative at $\lambda=L$.
It is the lowest-degree nonzero polynomial correction with this $C^2$ join.
The position of the lower endpoint and the cubic shape are constitutive
choices. Higher-order smooth corrections are possible.

For $\lambda>L$, variation gives $X=1/2$ and the original action exactly.
For $L/2<\lambda<L$, put $r=\lambda/L$. Then
\begin{align}
 X=e_\lambda&=\tfrac12-2(1-r)^2, & \tfrac12<r<1,\label{eq:X}\\
 e_{\lambda\lambda}&=\frac{4(1-r)}{L}>0, &
 P&=\lambda X-e-V.
\end{align}
The constrained endpoint gives a regular spacelike continuation. For $X<0$,
the maximizing auxiliary value is $\lambda=L/2$, hence
\begin{equation}
 P=\frac{L(\phi)}{2}X-K.
\end{equation}
This branch is a scalar with a positive field-dependent kinetic coefficient
and constant potential. At $X=0$ it joins continuously with the pressure branch.
The action is understood with the auxiliary inequality, or through its local
Hamiltonian. Its dust sector retains the usual potential-flow and caustic
limitations.

\subsection{Positive energy, causal characteristics and current}
For an observer with field time derivative $v$ and spatial gradient magnitude
$s$, the energy density is
\begin{equation}
 T_{00}=\frac{\lambda}{2}(v^2+s^2)+e+V.
\end{equation}
On the intermediate branch,
\begin{equation}
 \frac{e}{L}-\frac13=\frac{(2-r)(2r-1)^2}{6}\ge0.
\end{equation}
Thus $T_{00}\ge K$. The dust and spacelike branches obey the same lower
bound. The null contraction is $T_{\mu\nu}k^\mu k^\nu
=\lambda(k^\mu\nabla_\mu\phi)^2\ge0$.

Eliminating the auxiliary variable gives $P_X=\lambda$ and
$P_{XX}=1/e_{\lambda\lambda}$ on the regular pressure branch.
The local sound speed is therefore\cite{Garriga}
\begin{equation}\label{eq:sound}
 c_s^2=\frac{\lambda e_{\lambda\lambda}}
 {\lambda e_{\lambda\lambda}+2X}
 =\frac{4r(1-r)}{4r(1-r)+2X},\qquad 0\le c_s^2\le1.
\end{equation}
It tends to zero at the dust join and to one at $X=0$. These inequalities
control the local principal symbol. They do not prove global smooth
existence or prevent caustics and shocks.

For a regular timelike field, choose the future-directed flow and write
$v=\sqrt{2X}$, $j=\lambda v$ and $\theta=\nabla_\mu u^\mu$.
The scalar equation becomes
\begin{equation}\label{eq:current}
 \dot j+\theta j=-V_\phi F(r),\qquad
 F(r)=-1+6r^2-4r^3.
\end{equation}
The derivative of the moving threshold in $e$ is essential in this equation.
Direct factorization gives
\begin{align}
 F(r)&=(2r-1)(1+2r-2r^2),\\
 v^2&=(2r-1)(3-2r),\\
 F(r)&=A(r)v^2,\qquad
 A(r)=\frac{1+2r-2r^2}{3-2r}\in[\tfrac34,1].
\end{align}
Consequently,
\begin{equation}\label{eq:regular}
 \dot j=-\left[\theta+\frac{V_\phi A(r)j}{\lambda^2}\right]j.
\end{equation}
For regular flow and $K-V$ bounded away from zero, this fractional rate
remains continuous at $j=0$. Integrating-factor positivity now applies.
An initially positive future-directed current cannot cross zero in finite
regular evolution. With $V_\phi\ge0$, $\dot V=V_\phi v\ge0$.
The scalar slows relative to proper time as depletion approaches. The
unit-clock condition has been relinquished in this regime.

The constant minimum energy has a specific role. In a homogeneous element
at zero current, the scalar force is the derivative of the minimum total
energy. Requiring zero current to be an invariant state sets that derivative
to zero. A constant density cutoff instead leaves a nonzero force and allows
the clock to reverse. That unsuccessful alternative and its numerical
results are preserved in the research directory.

Covariant variation of~\eqref{eq:action} supplies the full stress tensor.
Its divergence vanishes on the field equations, with the baryon sector
included. Spatial variations of $V$ in the depleted phase are accompanied
by the pressure and momentum terms of the same action. A homogeneous vacuum
with an independently clipped drain is not being assumed there.

\subsection{An explicit choice of floor and inheritance of the fits}
Write the existing background as
\begin{equation}
 C(a)=B a^{-3}e^{-\eta_*F(a)},\qquad
 V(a)=\eta_* B\int^a a'^{-3}e^{-\eta_*F(a')}\,dF(a').
\end{equation}
Assume the future history is monotone with $F\le1$ and set today's scale
factor to one. Since $a^{-3}\le1$ in the future,
\begin{equation}
 V(\infty)-V(1)\le B[e^{-\eta_*F(1)}-e^{-\eta_*}].
\end{equation}
The demonstration selects $K$ as $V(1)$ plus this upper bound. This is a
specified effective choice using the inherited history and stock budget.
It is not a new fit to voids, and the microscopic theory has not selected it.

With densities in units $3H_{100}^2/(8\pi G)$, the selected value is
$K\simeq0.3369963$. The largest ratio $L/\bar C$ through today is
$0.0299937$. The homogeneous solution therefore lies inside the dust
branch with a substantial open neighborhood. The action and all its
variations there agree with the original action. The tested background and
linear perturbation equations, including their initial conditions, are
identical at these parameters.

This identity carries over the existing conditional linear scores:
$2227.80$ for Planck-lite TTTEEE+BAO+DES SN, $162.73$ for the separately
evaluated ACT primary-CMB summary, and $16.21$ for ACT lensing. These are
inherited results from the earlier implementation, with its massless-neutrino,
linear-lensing and other stated assumptions. No new likelihood was run.
Nonlinear structure, lensing and void statistics are not covered by this
identity, and overlapping likelihood scores must not be added.

\subsection{Numerical checks}
Five homogeneous curved FRW patches were evolved from $a_i=0.01$ with
initial contrasts $0,-0.1,-0.35,-0.6,-0.9$. Each patch conserves its own
baryons and radiation; curvature is fixed from its initial Friedmann
constraint. These patches test the local equations and are not compensated
void predictions. Every tested patch retained positive current and energy,
and its vacuum contribution remained monotone. Doubling the background grid
and tightening the integration gives a maximum normalized Friedmann residual
\MaxConstraintError. The largest change in a final reported dimensionless
state variable is \MaxPatchDifference.

\begin{figure}[ht]
\centering\includegraphics[width=\textwidth]{adaptive_boundary.png}
\caption{Selected constitutive completion in the homogeneous patch tests.
The auxiliary density approaches its positive endpoint and total dark energy
approaches $K$. Spatial interfaces and a compensating exterior are absent.}
\end{figure}

A separate periodic one-dimensional field calculation checks the regular
pressure branch at constant potential. At $r=3/4$, equation~\eqref{eq:sound}
predicts $c_s^2=1/2$. A perturbation of amplitude $10^{-4}$ gives a fitted
frequency differing from $1/\sqrt2$ by $3.7\times10^{-9}$ fractionally;
doubling the grid changes the frequency by $1.2\times10^{-12}$ in the chosen
units. Energy and mean current are conserved to floating-point precision in
this small-amplitude test. This check does not establish interface stability
or nonlinear global well-posedness.

\begin{figure}[ht]
\centering\includegraphics[width=.82\textwidth]{local_wave.png}
\caption{Regular local wave test. The field evolution and the sound speed
calculated from the action agree at the tested amplitude and resolutions.}
\end{figure}

\section{Conservative stock and the interaction still needed}
\subsection{A common action is available; equilibrium selection is separate}
A covariant positive-stock realization can use worldlines with
$m(\varphi)>0$ and a canonical receiver:
\begin{equation}\label{eq:receiver}
 S=\int\sqrt{-g}\left[\frac{M_{\rm Pl}^2}{2}\mathcal R
 -\tfrac12(\nabla\varphi)^2-U(\varphi)\right]+S_b[g]
 -\sum_i\int m(\varphi)\,ds_i.
\end{equation}
This is an effective coupled-scalar construction\cite{Amendola}. With
$\alpha=d\ln m/d\varphi$, its exchange equations are
\begin{equation}
 \nabla_\mu T_C^{\mu\nu}=-\alpha\rho_C\nabla^\nu\varphi,
 \qquad
 \nabla_\mu T_\varphi^{\mu\nu}=+\alpha\rho_C\nabla^\nu\varphi.
\end{equation}
The receiver carries the missing energy and momentum. Conserved particle
number and positive mass keep the particle stock positive. The continuum
kinetic description handles multiple streams. When the scalar exchange and
additional constitutive forces are negligible, it reduces to Vlasov--Poisson
evolution in the common baryon-plus-stock potential.

The new low-density action and the canonical worldline action are two
different effective completions. They have not been identified as one
multistream theory. A Hamiltonian director coupling can conserve total
momentum if its reaction forces and field momentum are included by variation
of a common translation-invariant energy. Its force and inertia terms still
need specification. Applying the full algebraic RAR enhancement in addition
to the gravitating stock counts the same excess twice.

Conversion by itself cannot provide the missing force. For a common
transport operator $\mathcal D$,
\begin{equation}
 \mathcal D f_u=-\Gamma f_u,\qquad
 \mathcal D f_R=+\Gamma f_u
 \quad\Longrightarrow\quad
 \mathcal D(f_u+f_R)=0.
\end{equation}
The total source and every total kinetic moment are unchanged by this
relabeling. A physical interaction must change the distribution if conversion
is to produce a different halo profile. Collisionless evolution alone admits
many equilibria. The circular-shell realization already given in O.3 can
support distinct positive spherical targets with equal total stock. For
example, Plummer stock profiles with scales $10$ and $40$ kpc and the same
mass give different forces at $20$ kpc. This is a direct nonuniqueness witness;
collective stability of each witness is a further question.

\subsection{The cosmological drain is insufficient to supply the tracking}
Consider an isolated, self-gravitating stock halo with fixed particle number,
slowly varying common particle mass $m$, homologous evolution, and conserved
canonical orbital actions. The scalings follow from
\begin{equation}
 J_{\rm orb}\sim mvr=\text{constant},\qquad
 v^2\sim\frac{GN m}{r}:
 \qquad r\propto m^{-3},\quad v\propto m^2.
\end{equation}
For fixed baryon mass, the acceleration scale inferred from $v^4/(GM_b)$
therefore scales as $m^8$. Between $z=2.3$ and today, the saved release history
gives $m_0/m_{2.3}=0.989398$. This produces an acceleration-scale ratio
$0.9183$, while $H_0/H(2.3)=0.2972$. Even a full unit of release gives only
$e^{-8\eta_*}=0.7875$ in this model. Nonadiabatic assembly, scalar gradients
and additional interactions are outside this restricted calculation. It
shows that the existing small mass drain is not an automatic explanation
of the proposed galactic tracking.

\subsection{A finite stock budget restricts the equilibrium domain}
For a pointlike baryonic mass, put $r_M=\sqrt{GM_b/a_0}$. The exponential
RAR requires
\begin{equation}
 \frac{M_R(<r)}{M_b}=\frac1{e^{r_M/r}-1}.
\end{equation}
A total available stock $kM_b$ therefore limits exact agreement to
\begin{equation}
 r\le r_{\rm cap}=\frac{r_M}{\ln(1+1/k)}.
\end{equation}
Here $k$ refers to the actual Lagrangian cohort and transport history.
It is not universally $5.1761$ per present baryon. The manuscript's Milky-Way
benchmark has about $14$ times its present baryonic mass in stock. The cap
principle is already present in the paper; its consequence is that an exact
unlimited RAR, a finite isolated budget and no imported source cannot all
hold. The outer transition and cohort assignment remain physical predictions.

\section{Response history and quantum noise}
\subsection{Passive lag has a definite sign}
Let $q(t)$ denote the equilibrium occupation of a fixed baryonic configuration.
For the first-order response
\begin{equation}
 \dot n=\gamma(q-n),\qquad \gamma>0,
\end{equation}
the difference $d=n-q$ obeys $\dot d=-\gamma d-\dot q$.
If $\dot q\le0$ and $d(t_i)\ge0$, integrating-factor positivity gives
$n(t)\ge q(t)$. The same conclusion holds for two positive relaxation poles
in cascade, initialized at or above the initial target. More generally a
positive normalized causal averaging kernel preserves this inequality.

For fixed $g_b$ and $a_0\propto H$ with decreasing $H$,
$q=[\exp\sqrt{g_b/a_0}-1]^{-1}$ decreases. Ordinary lag therefore raises
the force relative to instantaneous equilibrium. It cannot suppress the
high-redshift overprediction under these assumptions. A newly assembled
system may start below equilibrium, and a changing baryonic profile may
change the sign of $\dot q$. These cases require an independently specified
assembly and response history.

\begin{figure}[ht]
\centering\includegraphics[width=.9\textwidth]{response_history.png}
\caption{Fixed $g_b=2\times10^{-10}\,\mathrm{m\,s^{-2}}$, formation at $z=6$,
and two equal poles with the stated time constants. The initially equilibrated
histories remain above their instantaneous target. An initially empty
response can lie below it. These are illustrative histories, not galaxy fits.}
\end{figure}

\subsection{A noise--delay inequality for the existing two-pole model}
For the susceptibility already used in N.8,
\begin{equation}
 \chi^R=-\chi_0\frac{\gamma\Omega}{(\gamma-i\omega)(\Omega-i\omega)},
 \qquad
 D=-\operatorname{Im}\chi^R>0\quad(\omega>0),
\end{equation}
the normalized response has mean delay
$\bar\tau=\gamma^{-1}+\Omega^{-1}$.
Use the symmetrized quantum thermal relation\cite{Clerk} and the horizon
temperature $\beta\hbar=2\pi/H$:
\begin{equation}
 \mathcal F\equiv\frac{\sigma_n^2}{\chi_0/\beta}
 =\frac{\beta\hbar}{\pi\chi_0}\int_0^\infty
 D(\omega)\coth(\beta\hbar\omega/2)\,d\omega.
\end{equation}
Since $\coth x\ge1$, direct integration gives, for $\Omega>\gamma$,
\begin{equation}
 \mathcal F\ge\frac{2\gamma\Omega}{H(\Omega-\gamma)}
                 \ln\frac\Omega\gamma.
\end{equation}
Writing $r=\Omega/\gamma$, the inequality
$\ln r\ge2(r-1)/(r+1)$ yields
\begin{equation}\label{eq:noisedelay}
 \boxed{\mathcal F\,H\bar\tau\ge4.}
\end{equation}
The coincident-pole limit gives the same bound. The logarithmic inequality
and product bound are kernel-checked; numerical quadrature tests the thermal
integral across eighteen rate pairs. Equation~\eqref{eq:noisedelay} is specific
to this two-pole response family and the chosen temperature.

For an adiabatically changing target, the fractional occupation lag is
approximately $\bar\tau|\dot q/q|$. At $x=\sqrt{g_b/a_0}=1$ during matter
domination, three-percent occupation tracking requires
$H\bar\tau\lesssim0.0253$. The bound then gives $\mathcal F\gtrsim158$.
Using the existing independent-cell force formula, three-percent instantaneous
force rms requires $N_{\rm eff}\gtrsim6.47\times10^4$ in that example.
These are joint requirements on a proposed contact. Measurement windows and
spatial covariance change the observable noise. No noise exclusion follows
without those inputs.

Detailed balance does not select a rate. If $L$ has stationary distribution
$p$ and satisfies detailed balance, so does $sL$ for every $s>0$. The tests
include three such generators with identical stationary populations and
different finite-time responses. The physical interaction is needed to fix
the rate and the carrier's relaxation regime.

\section{Tensor propagation and the remaining matching freedom}
\subsection{What the minimal action establishes}
For the Einstein term in~\eqref{eq:action} or~\eqref{eq:receiver}, constant
$M_{\rm Pl}$ and homogeneous minimally coupled scalar matter give the tensor
quadratic action
\begin{equation}
 S_T^{(2)}=\frac{M_{\rm Pl}^2}{8}\int dt\,d^3x\,a^3
 \left[\dot h_{ij}\dot h_{ij}-a^{-2}\partial_kh_{ij}\partial_kh_{ij}\right].
\end{equation}
For propagation along the third axis, the symmetric transverse-traceless
matrix has the unique form
\begin{equation}
 h_{ij}=\begin{pmatrix}h_+&h_\times&0\\h_\times&-h_+&0\\0&0&0\end{pmatrix},
 \qquad h_{ij}h_{ij}=2(h_+^2+h_\times^2).
\end{equation}
There are two tensor kinetic coordinates, positive energy for
$M_{\rm Pl}^2>0$, and luminal principal propagation. The clock/receiver scalar
remains a material degree of freedom. The statement does not assert that
the whole theory contains only two degrees of freedom, nor that all possible
detector responses have been computed.

The minimal action is a sufficient effective choice. The static scalar
proofs do not force it. A covariant clock defines an extrinsic shear
$\sigma_{\mu\nu}$; an allowed additional term
$\frac{s}{2}\sigma_{\mu\nu}\sigma^{\mu\nu}$ vanishes on an exactly homogeneous
FRW background and in the corresponding static zero-shift limit. It changes
the tensor kinetic coefficient to $M_{\rm Pl}^2+s$, giving
\begin{equation}
 c_T^2=\frac{M_{\rm Pl}^2}{M_{\rm Pl}^2+s}.
\end{equation}
Excluding or suppressing such operators requires an action-level premise or
matching calculation. The importance of the tensor-speed operator map is
also explicit in scalar-tensor analyses\cite{Ezquiaga}. The failed
equilibration of the archived native tensor campaign remains unresolved.

\subsection{Why the identities do not select the physical coefficients}
The reciprocal relation permits the family
\begin{equation}
 \epsilon_\ell=\ell\frac{g}{(2\pi)^2},\qquad
 a_0=\epsilon_\ell cH,\qquad
 R_{\rm bare}=\epsilon_\ell^{-1},\qquad
 a_0R_{\rm bare}=cH,
\end{equation}
for every positive loading $\ell$. The two-phase symmetry and cancellation
identity survive while both individual predictions change. The stipulated
one-entropy loading chooses $\ell=1$. Selecting that loading physically
requires more than the reciprocal identity. This observation does not
invalidate the proven uniqueness inside a fully specified finite action.

There is an analogous background degeneracy in the canonical receiver.
For any smooth $0<W<C$, set
\begin{equation}
 D=C-W,\qquad \dot\varphi^2=W,\qquad U=V+W/2,
 \qquad m(\varphi(a))\propto a^3D(a).
\end{equation}
Then $D+\rho_\varphi=C+V$ and $p_\varphi=-V$ exactly.
The different choices of $W$ generally have different perturbations.
For the constant-fraction family $W=fC$, writing
$g=\eta_*dF/d\ln a$ and $\Omega_C=C/(3M_{\rm Pl}^2H^2)$ gives
\begin{equation}
 \alpha\equiv M_{\rm Pl}\frac{d\ln m}{d\varphi}
 =-\frac{g}{\sqrt{3f\Omega_C}},\qquad
 f(2\alpha^2)=\frac{2g^2}{3\Omega_C}.
\end{equation}
The maximum product on the saved history is $1.2752\times10^{-4}$.
The unscreened stock self-force enhancement is consequently about $12.75\%$
for $f=0.001$ and $2.55\%$ for $f=0.005$. These are conditional subhorizon
force factors, not measured total galactic forces. Scalar masses and spatial
response remain relevant. Sending $f$ to zero does not simultaneously remove
the receiver and its coupling in this family.

The routing and Farey proofs retain their conditional content. Symmetry,
equilibrium and background identities leave the physical loading, transition
matrix elements and some effective operators unselected. Historical
selection of formulas is an evidential issue that an additional theorem
cannot retroactively change.

\section{The remaining effective derivation}
The stock and carrier questions can be addressed by one specified physical
contact. At the effective level it must provide a Hamiltonian, the operators
coupling stock to the retained environment, and the environment's spatial and
temporal correlation functions. A weak-coupling realization would determine
transition rates from matrix elements and bath spectra. Detailed balance
would then fix their ratios; the matrix elements would set their scale.
The same interaction must account for mechanical energy, any protected
angular observables, and the response to baryonic assembly.

This is a narrower task than deriving every microscopic degree of freedom.
It still requires a physical choice beyond the supplied static and algebraic
premises. A useful candidate must recover the admissible RAR target, carry
the reaction forces, satisfy the stock budget, and predict its departures
before fitting each galaxy. The present work has not produced that operator.
It therefore does not resolve the high-redshift empirical tension, predict
cluster lens maps, or select a population of halo shapes.

For local release, equations~\eqref{eq:action}--\eqref{eq:regular} give a concrete
effective option with a stated physical cost. The next test is an
inhomogeneous calculation of interfaces and underdense structure, together
with a compatible multistream formulation. The selected floor and cubic can
be presented as a completion hypothesis. They should not be described as
already forced by the microscopic theory.

The tensor result closes the leading effective propagation calculation once
the minimal curvature action is chosen. The native geometric continuum,
its operator matching and its equilibrated numerical evidence remain separate
requirements. None of these open tasks is equivalent to a few additional
numerical certificates.

\section{Verification and reproduction}
The standalone supplement contains 42 kernel-checked theorems with only
\texttt{propext}, \texttt{Classical.choice} and \texttt{Quot.sound} in their
reported axiom dependencies. The checked statements cover energy and sound
speed inequalities, current-rate factorization, integrating-factor positivity,
passive lag, the logarithmic bound, receiver identities, matching freedom and
tensor-coordinate algebra. The full covariant variation and the spectral
integral are supplied as analytic derivations, with separate symbolic and
numerical checks.

The research directory contains the input hashes, a proof log, the source
files used for every run, both the first and refined patch results, wave
convergence results, and the first constant-cutoff attempt. The four current
papers and upstream Lean sources were preserved. No new observational
likelihood, full upstream rebuild, native tensor measurement or joint
galaxy/cluster simulation is claimed.

\begin{thebibliography}{9}
\bibitem{Garriga} J. Garriga and V. F. Mukhanov,
``Perturbations in k-inflation,'' Phys. Lett. B 458, 219 (1999),
\href{https://arxiv.org/abs/hep-th/9904176}{arXiv:hep-th/9904176}.
\bibitem{Amendola} L. Amendola, ``Coupled Quintessence,''
Phys. Rev. D 62, 043511 (2000),
\href{https://arxiv.org/abs/astro-ph/9908023}{arXiv:astro-ph/9908023}.
\bibitem{Clerk} A. A. Clerk et al., ``Introduction to Quantum Noise,
Measurement and Amplification,'' Rev. Mod. Phys. 82, 1155 (2010),
\href{https://arxiv.org/abs/0810.4729}{arXiv:0810.4729}.
\bibitem{Ezquiaga} J. M. Ezquiaga and M. Zumalac\'arregui,
``Dark Energy after GW170817: dead ends and the road ahead,''
Phys. Rev. Lett. 119, 251304 (2017),
\href{https://arxiv.org/abs/1710.05901}{arXiv:1710.05901}.
\end{thebibliography}
\end{document}
